\chapter{A Génese Dependente II: Surgimento Momentâneo}

\lettrine{N}{o livro de Ajahn Buddhadasa} sobre a Originação Dependente, ele
salienta que a sua abordagem tem incidido sobre o \emph{paṭiccasamuppāda} como
algo que se manifesta no momento presente, em vez de em termos de vidas
passadas, presentes e futuras. Quando contemplamos, quando praticamos,
percebemos que essa é a única forma como poderia ser. Isto porque estamos a
trabalhar com a própria mente. Mesmo quando consideramos o nascimento de um
corpo humano, não estamos a comentar o nascimento dos nossos próprios corpos,
mas a reconhecer mentalmente que esses corpos nasceram. Então, na reflexão,
observamos que a consciência mental surge e cessa. Assim, toda essa sequência da
Originação Dependente surge e cessa num instante. O surgimento e a cessação a
partir de \emph{avijjā} são momentâneos, não se trata de um tipo de
\emph{avijjā} permanente. Seria uma visão errada supor que tudo começou com
\emph{avijjā} e que, em algum momento no futuro, tudo cessaria.

\emph{Avijjā} significa, neste sentido, «não compreender as Quatro Nobres
Verdades». Quando há compreensão do Sofrimento, da Origem, da Cessação e do
Caminho, então as coisas já não são afetadas por \emph{avijjā}. Quando vemos com
\emph{vijjā,} então as perceções são realidade convencional, já não «eu» e
«meu». Por exemplo, quando há \emph{vijjā}, então posso dizer «Eu sou Ajahn
Sumedho» -- isso é uma realidade convencional, ainda uma perceção, mas já não é
vista a partir de \emph{avijjā,} é meramente uma convenção que usamos. Não há
nada mais do que isso. É como é.

Quando chegamos à cessação da ignorância, nesse momento, todo o resto da
sequência cessa. Não é como se uma coisa cessasse \emph{e}, \emph{em seguida},
outra cessasse. Quando há \emph{vijjā}, o sofrimento cessa. Em qualquer momento,
quando há verdadeira atenção plena e sabedoria, não há sofrimento. O sofrimento
cessou. Agora, quando contemplamos a cessação do desejo, a cessação do apego, há
a cessação do devir, a cessação do renascimento e do sofrimento. Quando as
coisas cessam, quando tudo cessa, então há paz, não é verdade? Há conhecimento,
serenidade, vazio, não-eu. Estas são as palavras, os conceitos que descrevem a
cessação.

Quando pratico desta forma, é muito difícil encontrar qualquer sofrimento.
Percebo que não há sofrimento, exceto num momento de descuido em que nos
deixamos levar por algo. Assim, devido ao descuido, à falta de atenção e ao
esquecimento, ficamos presos a padrões mentais habituais (kármicos). Mas quando
nós percebemos que fomos descuidados, podemos deixar isso cessar, podemos deixar
ir. Há o deixar ir, o permanecer no vazio. Já não existem os fortes impulsos
para agarrar; o fascínio e o encanto do mundo sensorial foram penetrados. Já não
há nada para agarrar. Ainda se pode experimentar e ver as coisas como elas são,
sem as agarrar. Não há ninguém a agarrar nada, mas ainda aqui pode haver sentir,
ver e ouvir, saborear e tocar. Já não é transformado numa pessoa\ldots{} «eu e o
que é meu».

Para mim, a perceção importante é precisamente o quão momentânea é a
consciência. A tendência é perceber a consciência como algo de longo prazo, como
estar acordado e consciente como um estado permanente do ser, em vez de um
momento.

E, no entanto, \emph{viññāṇa} é sempre descrita como um momento, um momento
fugaz, um instante. Assim, em vez de assumirmos que \emph{avijjā} é um processo
contínuo desde o nascimento até à morte, podemos ver que, a qualquer momento,
pode haver \emph{vijjā} e tudo simplesmente cessa. A cessação de toda essa massa
de sofrimento pode ser realizada. Desapareceu! Onde está?

Praticar desta forma é continuar a examinar as coisas para que tudo seja visto
exatamente como é. Tudo é apenas o que é no momento. Quando vemos isso, a beleza
é apenas beleza no momento. A fealdade é apenas isso no momento. Não há qualquer
tentativa de solidificar ou prolongar isso de forma alguma, porque as coisas são
apenas o que são. A pessoa torna-se cada vez mais consciente do informe ou
nebuloso como sendo apenas o que é, em vez de algo que é ignorado, descartado ou
mal interpretado.

O problema da perceção é que ela tende a limitar-nos a estar apenas conscientes
de certos pontos. Temos tendência a estar conscientes em certos pontos
designados e a mudança natural, o fluxo e a corrente não são realmente notados.
A pessoa está apenas consciente em A, B, C, D, E, F, G -- os pontos entre A e B
nunca são realmente notados porque a pessoa está apenas realmente consciente nos
pontos designados de perceção. É por isso que, quando a mente se abre com
\emph{vijjā} e se torna recetiva, o Dhamma revela-se; há uma espécie de
revelação. A mente vazia, num estado de admiração, permite que a verdade se
revele -- já não através da perceção. É aqui que reside a verdade inefável, as
palavras falham-nos e é impossível traduzi-la em perceções ou conceitos.

Talvez agora estejas a começar a apreciar a ênfase que o Buda colocou: «Eu
ensino o sofrimento e o fim do sofrimento. Eu ensino apenas duas coisas\ldots{}
há sofrimento e há o fim do sofrimento.» Se tiveres apenas essa compreensão do
sofrimento e perceberes o fim do sofrimento, então estás libertado da
ignorância. Se tentares especular sobre como isso é, poderias chamá-lo de
«Nibbana, a felicidade suprema» -- mas «felicidade suprema» também não é bem
isso, pois não? «Felicidade suprema» soa a ficar «pedrado», a flutuar no ar, a
alcançar o Nibbana e a flutuar até ao teto.

Mas o Caminho é um caminho de compreensão; atenção plena e compreensão. Então, o
Caminho Óctuplo é desenvolvimento, \emph{bhāvanā}: desenvolver esse caminho para
a compreensão correta. Cada vez mais percebemos o vazio, o não-eu, a liberdade
de não estarmos apegados a nada; o que afeta o que dizemos, o que fazemos e como
vivemos na sociedade em que estamos, aumentando a sensação de serenidade e
calma.

A palavra Nibbana é geralmente definida como «desapego aos cinco
\emph{khandhas}», o que significa deixar de sentir uma noção de «eu» em relação
ao corpo e à mente -- \emph{rūpa, vedanā, saññā, saṅkhāra, viññāṇa}.
Contemplamos os cinco \emph{khandhas} não mais com \emph{avijjā}, mas com
\emph{vijjā.} Percebemos que todos eles são impermanentes, insatisfatórios e
não-eu. Então, Nibbana é a compreensão do desapego, em que a visão do eu cessa.
O corpo continua a respirar, por isso não se dissolve no ar, mas a identidade
errada de que «eu sou o corpo» dissolve-se. A identidade errada com
\emph{vedanā, saññā, saṅkhāra} e \emph{viññāṇa} -- tudo isso cessa. O eu
dissolve-se, não se consegue encontrar ninguém. Não se consegue encontrar-se a
si mesmo porque se \emph{é} a si mesmo.

Na visão da Originação Dependente ocorrendo ao longo de três vidas, os cinco
\emph{khandhas} são vistos como uma forma permanente desde o nascimento. O
corpo: sensações, perceções, formações mentais e consciência são considerados
como sendo contínuos desde o nascimento. Mas isso é uma suposição que fazemos --
e o reflexo do surgimento momentâneo aponta para a própria mente. O corpo não é
uma pessoa de qualquer forma, não é «eu» nem «meu», nunca foi, nunca será. Há
apenas a \emph{perceção} dele como «eu» e «meu». A \emph{crença} de que nasci.

Tenho uma certidão de nascimento para provar que este corpo nasceu. Carregamos
certidões de nascimento nas nossas mentes -- carregamos toda a história, as
memórias e assim por diante das nossas vidas, dando-nos esta sensação de
continuidade de uma pessoa desde o nascimento até ao momento presente. Mas a
análise da perceção por si só mostra que a perceção surge e cessa. Esta perceção
de mim como uma personalidade permanente é apenas um momento. Surge e cessa. A
consciência, também, é apenas momentânea e transmite as qualidades atrativas,
repulsivas e neutras do reino condicionado. Quando se vê isso claramente, então
já não há interesse nesse apego e na busca da felicidade, tentando renascer na
felicidade ou na beleza, no prazer, na segurança ou na proteção. O renascimento
é um apego ao reino condicionado, por isso deixamos isso ir. Os cinco
\emph{khandhas} continuam a ser os cinco \emph{khandhas}; são vistos pelo que
são: impermanentes, insatisfatórios e sem um eu.

Portanto, é um reflexo da verdade de como as coisas são -- é muito direto, muito
claro. Do confuso, amorfo, nebuloso, inseguro, instável e incerto para o certo
-- seja o que for, já não estamos a escolher o que preferimos, estamos apenas a
observar que tudo o que surge cessa. À medida que percebes isto através da tua
prática, então grande parte da confusão é apenas confusão, é um \emph{dhamma.} A
confusão é apenas confusão no momento, não é permanente nem o eu. Assim, o que
antes era um problema ou algo que nos iludia transforma-se num \emph{dhamma.} A
transformação não se dá através da mudança da condição, mas através da mudança
da atitude, da ignorância para a clareza.

As pessoas dizem: «Tudo isto é muito bom, mas e quanto ao amor e à compaixão?» O
\emph{desejo} por tudo isso é o obstáculo, não é? O amor não é problema quando
não há ilusão; quando não há eu, não há nada que impeça, bloqueie ou evite o
amor. Mas enquanto houver a ilusão do eu, o amor é apenas uma ideia pela qual
ansiamos, mas com a qual nos sentimos sempre desapontados, porque o eu está a
atrapalhar. A visão do eu está sempre a cegar-nos, fazendo-nos esquecer e
imaginar que não existe amor. Sentimo-nos alienados, solitários e perdidos
porque parece não haver amor, por isso culpamos outra pessoa, ou culpamos a nós
próprios, talvez por não sermos dignos de amor, ou tornamo-nos cínicos.

Mas o Buda apontou para isto e perguntou qual era o verdadeiro problema? É a
ilusão de um eu. É o apego a essa perceção. Isso afeta a consciência, pelo que
estamos sempre a criar separações e insatisfação, e a identificar-nos com
aquilo que não somos nós próprios. Assim que nos libertamos dessa ilusão, o amor
está sempre presente. Apenas não o conseguimos ver nem desfrutar quando estamos
cegos pelos nossos desejos e medos.

À medida que compreende isto cada vez mais, a sua fé aumenta e surge a
disposição para abdicar de tudo. Há um verdadeiro entusiasmo, uma alegria em
estar com as coisas tal como são.
